\begin{figure}[htbp]
\centering
\begin{tikzpicture}[x=1mm,y=1mm,font=\sffamily\small]
  \node[text=brand] at (26,47) {Plane-wave basis at $\Gamma$};
  \node[text=brand] at (88,47) {Wavelength and detail};
  \fill[paleblue] (26,22) circle (18);
  \foreach \x in {-24,-18,...,24}
    \foreach \y in {-18,-12,...,18} {
      \pgfmathtruncatemacro{\inside}{\x*\x+\y*\y<=324}
      \ifnum\inside=1
        \fill[brand] ({26+\x},{22+\y}) circle (.75);
      \else
        \fill[muted!35] ({26+\x},{22+\y}) circle (.55);
      \fi
    }
  \draw[teal,line width=.8pt] (26,22) circle (18);
  \draw[->,draw=muted,line width=.4pt] (1,22) -- (53,22);
  \draw[->,draw=muted,line width=.4pt] (26,0) -- (26,43);
  \node[anchor=west,text=muted] at (51,19) {$G_x$};
  \node[anchor=west,text=muted] at (27,42) {$G_y$};
  \node[fill=paleblue,inner sep=1pt,text=brand] at (26,22) {$\Gamma$};
  \draw[draw=muted!40,line width=.4pt] (66,31) -- (112,31);
  \draw[draw=muted!40,line width=.4pt] (66,12) -- (112,12);
  \draw[teal,line width=.8pt,domain=0:44,samples=100]
    plot ({67+\x},{31+3*sin(\x*16.3636)});
  \draw[brand,line width=.8pt,domain=0:44,samples=140]
    plot ({67+\x},{12+3*sin(\x*49.0909)});
  \node[text=ink] at (89,39) {Longer wavelength};
  \node[text=ink] at (89,20) {Shorter wavelength};
  \node[text=ink] at (26,-7) {Higher Ecut: more components};
  \node[text=ink] at (89,-7) {Short waves need more kinetic energy};
\end{tikzpicture}
\caption{Plane-wave cutoff in two dimensions. Blue points are reciprocal components retained inside the cutoff circle; calculations use a three-dimensional set. The right panel relates wavelength to spatial variation.}
\label{fig:plane-wave-cutoff-concept}
\end{figure}
